\documentclass[11pt]{article}
\usepackage[margin=1in]{geometry}
\usepackage[T1]{fontenc}
\usepackage[utf8]{inputenc}
\usepackage{amsmath,amssymb,amsthm}
\usepackage{enumitem}

\title{ICPC World Finals 2021\\K. Take On Meme}
\author{}
\date{}

\begin{document}
\maketitle

\section*{Problem Summary}

Each leaf contains one point $(x,y)$. For an internal node with children $c_1,\dots,c_k$, we choose one
child as the winner and assign sign $+1$ to that child and sign $-1$ to all other children. The node then
produces
\[
    p_{winner} - \sum_{i \ne winner} p_i,
\]
where every $p_i$ is some point that the corresponding child subtree can produce.

At the root, we want the largest possible squared distance from the origin.

\section*{Key Observations}

\begin{itemize}[leftmargin=*]
    \item For any subtree, only the \emph{convex hull} of its achievable points matters. The final objective
    is to maximize Euclidean norm, and a point inside the convex hull is never better than an extreme point
    in the same direction.

    \item So the real primitive we need is:
    given a direction vector $d$, find the achievable point whose dot product with $d$ is maximum.

    \item This can be done recursively. Suppose an internal node has children with achievable convex sets
    $S_1,\dots,S_k$. For child $i$:
    \begin{itemize}[leftmargin=*]
        \item let $A_i$ be the point of $S_i$ maximizing $d \cdot p$;
        \item let $B_i$ be the point of $S_i$ minimizing $d \cdot p$.
    \end{itemize}
    If child $j$ is chosen as winner, the best achievable point is
    \[
        A_j - \sum_{i \ne j} B_i.
    \]
    Therefore we only need one maximizing point and one minimizing point from every child.

    \item To reconstruct the whole convex hull of the root, we use a support-oracle version of QuickHull:
    if we already know two hull vertices $a$ and $b$, then the only way another hull vertex can lie between
    them is on the left side of the directed segment $a \to b$. Querying the direction perpendicular to
    $a \to b$ gives the farthest achievable point on that side. If this query returns neither endpoint, we
    recurse.

    \item To avoid ambiguity when several points have the same dot product, we break ties
    lexicographically with a secondary direction. This guarantees that each support query returns a vertex,
    not an arbitrary interior point of a hull edge.
\end{itemize}

\section*{Algorithm}

\begin{enumerate}[leftmargin=*]
    \item For a query direction, store two vectors:
    \[
        (d_1, d_2),
    \]
    and compare points lexicographically by
    \[
        (d_1 \cdot p,\; d_2 \cdot p).
    \]

    \item Define a recursive function \texttt{extreme(u, $d_1$, $d_2$)} that returns the achievable point of
    subtree $u$ maximizing this lexicographic key.

    \item For a leaf, return its own point.

    \item For an internal node with children $c_i$:
    \begin{itemize}[leftmargin=*]
        \item compute
        \[
            A_i = \texttt{extreme}(c_i, d_1, d_2),
            \qquad
            B_i = \texttt{extreme}(c_i, -d_1, -d_2),
        \]
        so $B_i$ is the lexicographically smallest point in direction $d_1$;
        \item let
        \[
            T = -\sum_i B_i;
        \]
        \item if child $j$ wins, the best point is
        \[
            T + A_j + B_j.
        \]
        Choose the child $j$ that maximizes this expression lexicographically.
    \end{itemize}

    \item To enumerate root hull vertices, start from two opposite support queries:
    \begin{itemize}[leftmargin=*]
        \item one in direction $(1,0)$ with tie-break $(0,1)$,
        \item one in direction $(-1,0)$ with tie-break $(0,-1)$.
    \end{itemize}

    \item For every directed segment $a \to b$ found so far, query the normal direction
    \[
        d_1 = \operatorname{rot}(b-a),
    \]
    where $\operatorname{rot}(x,y)=(-y,x)$, and use $d_2=b-a$ as tie-break.
    If the returned point $c$ lies strictly to the left of $a \to b$, recurse on $(a,c)$ and $(c,b)$.
    Otherwise there is no additional hull vertex on that arc.

    \item After all recursive calls finish, evaluate $\|p\|^2$ for every discovered hull vertex and take the
    maximum.
\end{enumerate}

\section*{Correctness Proof}

We prove that the algorithm returns the correct answer.

\paragraph{Lemma 1.}
For a fixed direction pair $(d_1,d_2)$, \texttt{extreme(u, $d_1$, $d_2$)} returns a lexicographically
maximum achievable point of subtree $u$.

\paragraph{Proof.}
We use induction on the subtree.

For a leaf, the claim is trivial.

Now consider an internal node with children $c_1,\dots,c_k$. Fix one winning child $j$. To maximize
the resulting point lexicographically in direction pair $(d_1,d_2)$:
\begin{itemize}[leftmargin=*]
    \item the winner must contribute the lexicographically largest point of $c_j$, namely $A_j$;
    \item every losing child must contribute the lexicographically smallest point of its subtree, namely
    $B_i$, because these points are subtracted.
\end{itemize}
So the best outcome for winner $j$ is exactly
\[
    A_j - \sum_{i \ne j} B_i = T + A_j + B_j.
\]
Taking the best $j$ therefore yields the lexicographically maximum achievable point of the current node.
By the induction hypothesis, all $A_i$ and $B_i$ are computed correctly. \qed

\paragraph{Lemma 2.}
Every support query used by the hull-construction step returns a vertex of the root convex hull.

\paragraph{Proof.}
The primary direction chooses a supporting line of the convex hull. Among all points on that supporting
line, the secondary direction chooses one extreme endpoint of the corresponding face. Such an endpoint is
a hull vertex. By Lemma 1, the algorithm returns exactly this vertex. \qed

\paragraph{Lemma 3.}
Let $a$ and $b$ be two hull vertices. Querying the normal direction $\operatorname{rot}(b-a)$ finds a
new hull vertex on the arc from $a$ to $b$ if and only if such a vertex exists.

\paragraph{Proof.}
All points strictly on the counterclockwise arc from $a$ to $b$ lie on or to the left of the directed line
$a \to b$. The farthest hull point in the left normal direction is therefore the one farthest from the line
through $a$ and $b$ on that side. If this farthest point is still one of the endpoints, then no hull vertex
lies strictly on that side between them. Otherwise the returned point is a new hull vertex on that arc. \qed

\paragraph{Lemma 4.}
The recursive hull-construction step discovers every vertex of the root convex hull.

\paragraph{Proof.}
Start from two opposite hull vertices. They split the hull boundary into two arcs. By Lemma 3, whenever
an arc contains an undiscovered vertex, the recursive step finds one and splits the arc into two smaller
arcs. Whenever an arc contains no undiscovered vertex, recursion stops on that arc. Thus every arc is
eventually resolved, and every hull vertex is discovered. \qed

\paragraph{Theorem.}
The algorithm outputs the correct answer for every valid input.

\paragraph{Proof.}
By Lemma 4, the algorithm enumerates all vertices of the root convex hull. Since the squared norm is a
convex function, its maximum over a convex polygon is attained at some vertex. Therefore taking the
largest $\|p\|^2$ among all discovered hull vertices gives the true optimum. \qed

\section*{Complexity Analysis}

Let $N$ be the number of tree nodes and let $H$ be the number of vertices of the root convex hull.

One call to \texttt{extreme} visits every node once, so it costs $O(N)$ time. The hull construction makes
$O(H)$ support queries, hence the total time is
\[
    O(NH).
\]
This is fast enough because the tree has height at most $10$ and the coordinate bounds keep the hull size
small in practice.

The memory usage is $O(N)$ for the tree and the recursion stack.

\section*{Implementation Notes}

\begin{itemize}[leftmargin=*]
    \item All coordinates fit in signed $64$-bit integers. Every leaf contributes with coefficient $\pm 1$, so
    absolute coordinates stay within about $10^7$.
    \item The recursion on hull arcs uses the sign of the cross product to detect whether a new point lies
    strictly to the left of the current segment.
    \item The implementation stores discovered hull vertices in a set, because different support directions
    can legitimately return the same vertex.
\end{itemize}

\end{document}
